\documentclass[10pt,a4paper]{amsart}
\usepackage[margin=19mm]{geometry}
\usepackage{amsmath,amssymb,mathtools}
\usepackage[scr=rsfs,cal=euler]{mathalfa}
\usepackage{xfrac}
\usepackage[unicode,hidelinks,bookmarks=false]{hyperref}
\newcommand{\ie}{i.e.\ }
\newcommand{\cf}{cf.\ }
\newcommand{\resp}{respectively\ }
\newcommand{\Subsection}{\subsection}
\providecommand{\nobreakdash}{\nobreak\hbox{-}}
\setlength{\emergencystretch}{2em}
\multlinegap0pt
\usepackage[shortlabels]{enumitem}
\usepackage{amssymb}
\usepackage{tikz}
\usepackage{xcolor}
\usepackage{float}
\usepackage{algorithm}\usepackage{algpseudocode}
\newtheorem{lemma}{Lemma}
\newtheorem{theorem}{Theorem}
\usepackage{cleveref}
\crefname{lemma}{Lemma}{Lemmas}
\crefname{theorem}{Theorem}{Theorems}
  \newcommand{\defeq}{\coloneq}
  \renewcommand{\phi}{\varphi}
\title{\textsc{ Backpropagation from KL Projections: Differential and Exact I-Projection Correspondences }}
\author{Manish Krishan Lal}
\date{Source: arXiv:2512.24335v2 (CC BY 4.0)}
\begin{document}
\maketitle

\noindent\emph{Source and licence.} \textsc{ Backpropagation from KL Projections: Differential and Exact I-Projection Correspondences }. Version arXiv:2512.24335v2 (CC BY 4.0), distributed by arXiv under the Creative Commons Attribution 4.0 International licence, http://creativecommons.org/licenses/by/4.0/, as recorded in the arXiv licence metadata for this exact version and shown on the abstract page https://arxiv.org/abs/2512.24335v2. Full licence notice: \texttt{licenses/arxiv-2512.24335-CC-BY-4.0.txt}.\par\smallskip

\noindent\textit{Editorial scope.} Counted unit: the article's theorem thm:caseA (source0.tex lines 1570--1580). The article's complete original proof of it is delivered with it (source0.tex lines 1582--1624, one \texttt{proof} environment of 43 lines). No mathematics is written, repaired or re-derived: the statement and the proof are the article's own lines, in the article's own notation, and no retained line is substituted or re-flowed.  Retained uncounted as the source-local context the proof needs: lines 1526--1535.  Nothing was omitted from any retained span, so the package carries no omission record.  No retained line contains a citation, so the wrapper carries no bibliography and no bibliography entry.  No retained line contains a figure environment, an image-inclusion command or drawing code, so no image is delivered and the package declares no asset dependency.  Editorial formatting only, in addition: an AMS \texttt{amsart} preamble that restates the article's own declarations and macros used by the retained lines, namely the environments \texttt{lemma}, \texttt{theorem} on separate counters, together with the standard packages those lines need.  The retained environments print in this document as Lemma 1, Theorem 1 in order of appearance, whereas the article prints the same items with the numbering of its own full document.  The complete native source of the article as distributed in the arXiv e-print for this exact version is bundled unchanged as \texttt{sources/191881/source0.tex}.
\begin{lemma}[Delta-factor chain rule]
\label{lem:delta-chain}
Under these assumptions,
\[
\left.\frac{\partial}{\partial x_j}\log m_{g\to x_j}(x_j)\right|_{x_j^*}
\;=\;
\left.\frac{\partial}{\partial y}\log m_{y\to g}(y)\right|_{y=y^*}\cdot
\left.\frac{\partial y}{\partial x_j}\right|_{x^*}.
\]
\end{lemma}

\begin{theorem}[Backpropagation as log-derivatives of downward BP messages]
\label{thm:caseA}
With the setup above, $s(z)=\phi'(z^*)/\phi(z^*)$, and for any variable $x$,
\[
\left.\nabla_x \log b_\downarrow(x)\right|_{x^*}
\;=\;
\frac{\phi'(z^*)}{\phi(z^*)}\ \left.\nabla_x z\right|_{x^*}.
\]
Equivalently, for $\phi(z)=e^{\alpha z}$ one has
$\nabla_x \log b_\downarrow(x)=\alpha\,\nabla_x z$ at the forward point $x^*$.
\end{theorem}

\begin{proof}
Let
\[
c\;\defeq\;\frac{\phi'(z^*)}{\phi(z^*)}.
\]
Then $s(z)=c$ by the choice of the seed message $m_{\phi\to z}=\phi$. At each deterministic factor $g$ corresponding to $y=\psi(\cdot)$,
\cref{lem:delta-chain} yields
\[
\partial_x\log m_{g\to x}
\;=\;
\partial_y\log m_{y\to g}\cdot \partial_x y
\quad\text{at the forward point}.
\]
Since $b_\downarrow(x)$ is the product of all factor-to-variable messages into $x$,
\[
s(x)
\;=\;
\partial_x\log b_\downarrow(x)
\;=\;
\sum_{y\in \mathrm{ch}(x)} s(y)\,\partial_x y,
\]
where $\mathrm{ch}(x)$ denotes the children of $x$ in the computation DAG. Thus the family $x\mapsto s(x)$ satisfies the same backward recursion as the family $x\mapsto c\,\nabla_x z$. Indeed, by the ordinary chain rule,
\[
c\,\nabla_x z
\;=\;
\sum_{y\in\mathrm{ch}(x)} c\,\nabla_y z\,\partial_x y.
\]
To conclude, choose a topological ordering $v_1,\dots,v_N$ of the DAG such that every edge points from a lower index to a higher index, and let $v_N=z$. We prove by backward induction on this order that
\[
s(v_k)=c\,\nabla_{v_k}z\qquad\text{for }k=N,N-1,\dots,1.
\]
For $k=N$ this is exactly the identity $s(z)=c=c\,\nabla_z z$. Suppose now that the identity holds for every child $y\in\mathrm{ch}(x)$ of a node $x=v_k$. Then
\[
s(x)
=
\sum_{y\in\mathrm{ch}(x)} s(y)\,\partial_x y
=
\sum_{y\in\mathrm{ch}(x)} c\,\nabla_y z\,\partial_x y
=
c\,\nabla_x z,
\]
where the last step is the chain rule above. This closes the induction and proves the theorem.
\end{proof}

\end{document}
